\section{Exemple de plan}

\section*{\centering 1. Introduction ($\approx 1$ min)}
\begin{enumerate}[label=\arabic*.]
    \item \textbf{Accroche Impactante :} L'IA de recrutement d'Amazon a été \textbf{abandonnée} car elle a appris à pénaliser systématiquement les CV mentionnant des expériences féminines. Le numérique, loin d'être neutre, reproduit nos biais.
    \item \textbf{Problématique :} \emph{Comment lutter contre les biais algorithmiques pour garantir l'équité des systèmes d'Intelligence Artificielle ?}
    \item \textbf{Définition (Par Analogie) :} L'IA est une \textbf{recette de cuisine} qui s'auto-écrit à partir de ses \textbf{ingrédients} (les données). Le \textbf{biais} est l'erreur systématique qui rend le plat injuste pour certains convives.
    \item \textbf{Motivation (NSI/Éthique) :} Mon intérêt s'est affirmé en NSI : j'ai compris que nos \textbf{choix algorithmiques} (paramètres, modèles) ont un \textbf{impact éthique} décisif sur la société.
    \item \textbf{Annonce du Plan :} Nous analyserons d'abord les \textbf{causes profondes} (données et modélisation), puis les \textbf{solutions} (correction technique et cadre légal/éthique).
\end{enumerate}

\hrule

\section*{\centering 2. Axe de Développement 1 : Les causes profondes des biais ($\approx 2$ min)}
\subsection*{Intitulé : Le problème est dans les données et la modélisation, pas le code seul.}
\begin{itemize}[label=$\diamond$]
    \item \textbf{Biais d'échantillonnage (Le mauvais apprentissage) :}
    \begin{itemize}[label=\textendash]
        \item Si une classe est sous-représentée dans la \textbf{base de données} (NSI : Données structurées), l'IA se trompe sur elle.
        \item \textbf{Exemple Illlustratif :} Le modèle classe toujours le chien \emph{Husky} uniquement par la présence de \textbf{neige} en arrière-plan dans la base d'images. Il apprend le contexte, pas l'animal.
    \end{itemize}
    \item \textbf{Biais Historique (Le reflet du passé) :}
    \begin{itemize}[label=\textendash]
        \item Les jeux de données sont un miroir des discriminations passées (NSI : Pertinence des données).
        \item \textbf{Explication (Amazon) :} L'algorithme a appris à pénaliser les femmes car l'historique d'embauche était biaisé, il a \textbf{reproduit} la discrimination.
    \end{itemize}
    \item \textbf{Biais de Modélisation (Les mauvaises questions) :}
    \begin{itemize}[label=\textendash]
        \item Le développeur choisit ce que l'IA doit maximiser. (NSI : Algorithmique, Modélisation).
        \item \textbf{Exemple : } Le bot \emph{Tay} de Microsoft est devenu raciste et sexiste car il a été programmé pour \textbf{maximiser l'engagement}, en apprenant des utilisateurs extrêmes.
    \end{itemize}
    \item \textbf{L'Opacité de la "Boîte Noire" :}
    \begin{itemize}[label=\textendash]
        \item La complexité des \textbf{réseaux de neurones} empêche de retracer le \emph{pourquoi} d'une décision. (NSI : Langages et programmation).
        \item \textbf{Analogie : } C'est un \textbf{oracle} qui donne des réponses sans expliquer ses calculs.
    \end{itemize}
\end{itemize}

\hrule

\section*{\centering 3. Axe de Développement 2 : Les solutions : Audits et Réglementation ($\approx 1$ min)}
\subsection*{Intitulé : Vers une IA responsable : l'ingénierie de l'équité et le cadre légal.}
\begin{itemize}[label=$\diamond$]
    \item \textbf{L'Ingénierie des données \emph{a priori} :}
    \begin{itemize}[label=\textendash]
        \item \textbf{Action :} Il faut corriger le déséquilibre par l'\textbf{oversampling} (on réplique les données des minorités) ou l'\textbf{undersampling} (on retire les données des majorités).
        \item \textbf{Méthode (NSI, Bases de données) :} On utilise des \textbf{requêtes SQL} (\texttt{COUNT}, \texttt{GROUP BY}) pour analyser la distribution des données avant toute chose.
    \end{itemize}
    \item \textbf{L'Explicabilité (\emph{XAI}) des modèles :}
    \begin{itemize}[label=\textendash]
        \item \textbf{Objectif :} Développer des méthodes pour que l'IA \textbf{justifie ses décisions} ("J'ai rejeté ce prêt car le facteur X est trop bas").
        \item \textbf{Analogie : } L'IA doit donner la \textbf{liste des ingrédients} qui ont mené au plat final.
    \end{itemize}
    \item \textbf{Robustesse et Sécurité face aux Attaques Adversariales :}
    \begin{itemize}[label=\textendash]
        \item \textbf{Problème :} Lutter contre les \textbf{attaques par empoisonnement} (corruption intentionnelle des données).
        \item \textbf{Exemple Concret (NSI, Réseaux) :} De simples \textbf{stickers} sur un panneau Stop peuvent le faire identifier comme un panneau de "limite de vitesse" par l'IA d'une voiture autonome. L'algorithme se focalise sur les pixels du sticker.
    \end{itemize}
    \item \textbf{Le Cadre Légal (Obligation d'explication) :}
    \begin{itemize}[label=\textendash]
        \item Le \textbf{RGPD} (Règlement Général sur la Protection des Données) oblige les entreprises à la \textbf{transparence} et à la \textbf{justification} des décisions automatisées. (NSI : Impact sociétal et éthique).
    \end{itemize}
\end{itemize}

\hrule

\section*{\centering 4. Conclusion ($\approx 1$ min)}
\begin{enumerate}[label=\arabic*.]
    \item \textbf{Réponse Synthétique :} La lutte contre les biais est une \textbf{boucle vertueuse} : elle exige de corriger la \textbf{qualité des données} (technique NSI) et d'imposer un \textbf{cadre éthique fort} pour exiger l'\textbf{explicabilité}.
    \item \textbf{Lien Personnel :} Ce sujet m'a montré qu'en tant que futur citoyen/professionnel, mon rôle est d'exiger une \textbf{transparence algorithmique} et une \textbf{responsabilité du code} pour un numérique plus juste.
    \item \textbf{Ouverture (Rationalisée) :} Si l'IA est plus juste, le danger reste la \textbf{Cybersécurité}. La prochaine étape, c'est de garantir la \textbf{confiance} dans ces systèmes de plus en plus autonomes.
    \item \textbf{Question de Relance (Stratégique) :} Comment la \textbf{cryptologie} peut-elle garantir à la fois l'anonymat des données pour le citoyen et la traçabilité des décisions d'une IA pour les autorités de contrôle ?
\end{enumerate}

\newpage

\renewcommand{\arraystretch}{4} 

\begin{table}[h!]
\centering
\begin{tabular}{p{0.5\linewidth}|p{0.5\linewidth}} 
\hline
\multicolumn{1}{c|}{\textbf{POINTS FORTS}} & \multicolumn{1}{c}{\textbf{POINTS FAIBLES}} \\ 
\hline
& \\ % Ligne 1
& \\ % Ligne 2
& \\ % Ligne 3
& \\ % Ligne 4
& \\ % Ligne 5
& \\ % Ligne 6
& \\ % Ligne 7
& \\ % Ligne 8
& \\ % Ligne 9
& \\ % Ligne 10
& \\ % Ligne 11
& \\ % Ligne 12
& \\ % Ligne 9
& \\ % Ligne 9
\hline
\end{tabular}
\end{table}

\newpage

\begin{table}[h!]
\centering
\begin{tabular}{p{0.5\linewidth}|p{0.5\linewidth}} 
\hline
\multicolumn{1}{c|}{\textbf{POINTS FORTS}} & \multicolumn{1}{c}{\textbf{POINTS FAIBLES}} \\ 
\hline
& \\ % Ligne 1
& \\ % Ligne 2
& \\ % Ligne 3
& \\ % Ligne 4
& \\ % Ligne 5
& \\ % Ligne 6
& \\ % Ligne 7
& \\ % Ligne 8
& \\ % Ligne 9
& \\ % Ligne 10
& \\ % Ligne 11
& \\ % Ligne 12
& \\ % Ligne 9
& \\ % Ligne 9
\hline
\end{tabular}
\end{table}